% 图 1-1 设计与验证流程（坐标单位 cm；标注置于框内）
\begin{tikzpicture}[font=\footnotesize,
  s/.style={draw=softline, rounded corners=1.5pt, fill=macblue!55, align=center, inner sep=1.8pt,
            font=\scriptsize, minimum width=3.55cm, minimum height=1.25cm, text width=3.3cm},
  ar/.style={-{Stealth[length=1.4mm]}, thick}]

\node[s, fill=maclav] (a) at (-5.85,0) {\textbf{指标定义与架构划分}\\官方模板五部分成文；\\指标逐项可回溯};
\node[s, fill=macblue] (b) at (-1.95,0) {\textbf{RTL 设计}\\CPU / Ascon / PYNQ；\\逐模块配 testbench};
\node[s, fill=macgreen] (c) at (1.95,0) {\textbf{单元与子系统仿真}\\Icarus / XSim；\\位精确比对（L0/L1）};
\node[s, fill=macgreen] (d) at (5.85,0) {\textbf{SoC 集成与回归}\\smoke 6 PASS；\\regression 15 PASS（L2）};

\node[s, fill=macamber] (e) at (5.85,-1.95) {\textbf{综合}\\Vivado 2025.2；\\xc7z020clg400-1};
\node[s, fill=macamber] (f) at (1.95,-1.95) {\textbf{实现与策略扫掠}\\33 组并行；\\回收 0.824 ns};
\node[s, fill=macrose] (g) at (-1.95,-1.95) {\textbf{归因改造}\\复位 / 乘法 / Cache；\\156 个失败端点};
\node[s, fill=macamber] (h) at (-5.85,-1.95) {\textbf{当前 PYNQ 实现}\\WNS $+0.034$ ns；\\DRC 0 Error / 0 Critical};

\node[s, fill=macpeach] (i) at (-5.85,-3.9) {\textbf{当前板级 KAT}\\三模式各 289/289\\错误 Tag 各 289/289};
\node[s, fill=macpeach] (j) at (-1.95,-3.9) {\textbf{当前加密周期}\\SW/HW/ISE\\23342 / 1439 / 3995};
\node[s, fill=macrose] (k) at (1.95,-3.9) {\textbf{错误恢复}\\9/9 checks；7 次 reset\\每模式 50 次往返};
\node[s, fill=macgray] (l) at (5.85,-3.9) {\textbf{证据与复现}\\发布哈希与原始样本\\clean-clone 待验收};

\draw[ar] (a) -- (b);
\draw[ar] (b) -- (c);
\draw[ar] (c) -- (d);
\draw[ar] (d.south) -- (e.north);
\draw[ar] (e) -- (f);
\draw[ar] (f) -- (g);
\draw[ar] (g) -- (h);
\draw[ar] (h.south) -- (i.north);
\draw[ar] (i) -- (j);
\draw[ar] (j) -- (k);
\draw[ar] (k) -- (l);
% 反馈：归因改造不收敛时回到实现与策略扫掠；文字放在弧线上方的空白处
\draw[ar, dashed] (g.north) to[bend left=8] node[midway, above=0pt, font=\scriptsize, inner sep=1.5pt] {不收敛则回改 RTL} (f.north);
\end{tikzpicture}
